\documentclass[english]{article}

\usepackage{ae,aecompl}
\usepackage[T1]{fontenc}
\usepackage[latin9]{inputenc}
\usepackage{textcomp}
\usepackage{graphicx}
\usepackage{eurosym}
\usepackage{hyperref}
\usepackage{float}
\usepackage{fancyhdr}
\pagestyle{fancy}
\lhead{PowerEnJoy RASD}

\newcommand{\carsharing}{\textit {car sharing }}
\newcommand{\powerenjoy}{\textit{PowerEnJoy }}
\newcommand{\registereduser}{\textit {registered user }}
\newcommand{\staff}{\textit{staff }}
\newcommand{\service}{\textit{service }}
\newcommand{\safearea}{\textit{safe area }}
\newcommand{\powergrid}{\textit{Power Grid }}
\newcommand{\resevation}{\textit{reservation }}
\newcommand{\stand}{\textit{stand }}
\usepackage{tabto}

\makeatletter
\usepackage{babel}


\makeatother

\usepackage{babel}
\begin{document}
\begin{figure}
	\centering
	\includegraphics[scale=0.5]{logo.pdf} 
\end{figure}


\title{PowerEnJoy\\
 Requirement Analysis and Specification Document\\
}

\date{A.A 2016/2017}

\author{Erba Alessandro\\
 Leveni Filippo\\
 Lodi Luca}

\maketitle
\pagebreak{}

\tableofcontents{} \pagebreak{}

\section{Introduction}
	\subsection{Purpose }
	
		The purpose of our team is to project PowerEnjoy, which is  a platform that will be used to manage a car sharing service 
  in the city area.\\
  The platform will allow clients to find and reserve an electric car using mobile app from their self phone, and it will also    manage the communication of the emergency situations to the proper staff.\\
  Furthermore, a particular PowerEnjoy feature is the application of certain discounts to its users.\\
  \tab \tab This document aims to describe the high-level functionalities that will be offered by PowerEnjoy service. The    RASD is intended to be viewed by the stakeholders, that will evaluate the correctness of any assumption and decision
  in this document.
  
 \subsection{Scope and Problem Description}
 
  The government of a large city would like to introduce an ecological way to travel within the city area. In particular,     uniform distribution of cars in the city is a significant target.\\
  \tab \tab The cars are electric as anticipated and inside them it is available an integrated interface. Through this interface   is possible to communicate with the system (e.g. to call emergency staff).\\
  \tab \tab Users can choose and reserve a car through a mobile app. The system answers to the request by showing the    time left to the expiration of the reservation.\\
  The staff is divided in three categories:\\
   - Field staff\\
   - Emergency staff\\
   - Management staff\\
  \tab \tab Field staff use the mobile application to inform the system about their availability and to confirm that they are    going to take care of a certain call (e.g. retrieve a broken car, move a vehicle with a low battery or to ensure the uniform    distribution in the city, etc...).\\
  \tab \tab Emergency staff use the mobile application to answer customers calls and to take care of special situations (e.g.    no money in the costumer's credit card, car accident, etc...).\\
  \tab \tab Management staff use the mobile application to handle the service (e.g. set the safe areas, fares, discounts,     etc...).\\
  \tab \tab The system guarantees a uniform distribution of cars in the city. In particular, the city is divided in       circumferences and sectors. The system automatically computes the distribution of available cars both in the areas     between the circumferences and in the sectors, based on the GPS information it receives from each car. Uniformity is     guaranteed by ensuring the same amount of available cars for each area between the circumferences and for each sector.   \\
  \tab \tab A user can take advantage of the "Money saving" option selecting it on the integrated interface in the car. After    that he/she can input his/her final destination and the system provides information about the station where to leave the    car to get a discount. This destination is determined to ensure the uniform distribution of cars and depends both     on the destination of the user and on the availability of power plugs at the selected station.
	
	\subsection{Goals}
	\subsection{Domain Assumptions }
	\subsection{Glossary}
		\subsubsection{Car Sharing}
			\carsharing is a model of car rental where people rent cars for short periods of time, often by the hour.
		\subsubsection{PowerEnJoy}
			\powerenjoy is the \carsharing brand object of this document. 
		\subsubsection{Registered User}
			A \registereduser is a person subscribed to \powerenjoy which can access to the \carsharing through his smart-phone.
		\subsubsection{Staff}
			\staff is the set of people that are \registereduser but can perform special operations. They are divided in three different categories.
			\begin{itemize}
				\item {Field Staff}
				\begin{itemize}
				\item They own a passepartout for cars that has all kind of issues, for example they can menage cars with low battery.
				\end{itemize}
				\item{Emergency Staff}
				\begin{itemize}
				\item They manage users issues. They are available to communicate with users.
				\end{itemize}
				\item{Management Staff}
				\begin{itemize}
				\item They manage the system, for example they can change system parameters. They can add new \safearea, modify \carsharing prices.
				\end{itemize}
			\end{itemize}
		\subsubsection{Service}
			The \service is the group of actions that a \registereduser can fulfill trough the mobile.	
		\subsubsection {Safe Area}
			A \safearea is a  geographical place\footnote{defined by a set of GPS positions} on the map in which parking is allowed. \safearea are saved into the system. A \registereduser can end a \carsharing parking his car only in safe areas.
	\subsubsection{Power Grid}
		A \powergrid is the station that charge the battery of an electric car. All \powergrid are displayed in the map on-board \powerenjoy cars. Each \powergrid is in a \safearea. A \registereduser who plugs his car into a \powergrid receives the discount expected within \powerenjoy rules.
	\subsubsection{Reservation}
		A \resevation is the possibility that a \registereduser has to book a car at the latest for one hour.
	\subsubsection{Stand}
		A \stand is the possibility that a \registereduser has to pause his \carsharing. During a \stand the \registereduser can leave the car, than he can re-unlock it through his samrt-phone 

\subsection{Further Developments}
	We think that one of the most important things is to release a Service that can be extended and improved. It's trivial that a service like this must be able to be change and be updated.
	For this purpose we will develop a system open to new extensions. \\
	Here we provide a series of possible extensions that could be necessary in the future.
	E.g.:
	\begin{itemize}
		\item New kind of discount policies can be introduced to the system, for example discounts for new users. 
		\item New system policies may be necessary, for example new \carsharing conditions.
		\item New kind of vehicles can be introduced, for example electric scooter and electric bikes.
		\item Support to new payment methods, for example Paypal\textregistered, NFC\footnote{Near Field Communication; NFC devices are used in contactless payment systems, similar to those used in credit cards and electronic ticket smartcards and allow mobile payment to replace/supplement these systems. \href{https://en.wikipedia.org/wiki/Near_field_communication}{NFC on Wikipedia}}.
		\item Develop new services in order to increase the user experience, for example it may be useful to allow communication among users in order to share a car ride and save money.
	\end{itemize}
	\subsection{Used Tools}
	\begin{itemize}
		\item \LaTeX\\
		\item GitHub\\
	\end{itemize}
\section{Specific Requirements}
	\subsection{Functional Requirements}
	\subsection{Non Functional Requirements}

\section{Scenarios Identifying}
\section{UML Models And Use Cases}
	\subsection{Use Cases Diagram}
	\subsection{Actors Identifying}
	\subsection{Use Cases}

	\subsection{Class Diagram}
\section{External Interfaces}
\section{Alloy Model}
\section{Hours of Work}

\end{document}